\section{Problems for Stage I}
\label{problems:stage-I}

\begin{enumerate}[label=\textbf{I.\arabic*},leftmargin=*]
  \item Show that the Born rule for an effect $0\leq F\leq\identity$ gives a
  number in $[0,1]$.  Construct a two-outcome POVM that is not a projection-valued measure and give it a finite-resolution detector interpretation.

  \item For the momentum operator on $(0,L)$, verify by integration by parts
  that the boundary condition $\psi(L)=\rme^{\rmi\alpha}\psi(0)$ makes the
  operator self-adjoint.  Derive Eq.~\eqref{eq:momentum-extension-spectrum}.

  \item Let $e_n$ be an orthonormal sequence.  Prove that $e_n$ converges
  weakly to zero but not in norm.  Interpret $e_n(x)=\phi(x-n)$ on the line as
  probability escaping every fixed compact detector.

  \item Starting from $\rmd^3p=p^2\dd p\rmd\Omega$ and
  $E=p^2/(2\mu)$, derive the density factor in
  Eq.~\eqref{eq:free-energy-unitary} and verify unitarity explicitly.

  \item Prove that a bounded operator commuting with every multiplication
  operator $f(E)$ in a multiplicity-one direct integral is itself
  multiplication by an essentially bounded function.  State the change when
  the fiber has dimension $m(E)>1$.

  \item Calculate the spectral measure of a normalized vector for a
  two-level Hamiltonian.  Repeat for a wave packet of the one-dimensional free
  Hamiltonian and identify the Jacobian between momentum and energy density.

  \item Use the resolvent identity to prove analyticity of
  $\Resolvent(z)=(z-H)^{-1}$ in operator norm on the resolvent set.  Explain why this
  statement alone cannot continue $\Resolvent$ through the continuous spectrum.

  \item Derive the Feshbach operator
  Eq.~\eqref{eq:feshbach-preview}.  Differentiate its eigenvalue equation with
  respect to $E$ and identify the origin of
  $1-\partial_E\vsub{H}{eff}(E)$ in pole normalization.

  \item For a Lorentzian density cut off below a threshold, evaluate the
  contour contribution to Eq.~\eqref{eq:survival-spectral-fourier}.  Separate
  the pole term from the threshold term and determine their qualitative time
  dependence.

  \item Give one example each of norm, strong-operator, weak-operator, and
  strong-resolvent convergence.  For every example, name a physical statement
  that would be false if the topology were silently strengthened.

  \item Let $P_N\to\identity$ strongly and let $K$ be compact.  Prove
  $\|(\identity-P_N)K\|\to0$ by covering the image of the unit ball with
  finitely many $\varepsilon$ balls.  Give a counterexample when $K$ is the
  identity.

  \item For the three-dimensional off-shell free Green kernel and a bounded
  compactly supported potential, verify the finiteness of
  Eq.~\eqref{eq:bs-hs-integral}.  Treat separately the diagonal singularity
  and large separations.

  \item Prove the Birman--Schwinger equivalence
  Eq.~\eqref{eq:birman-schwinger-equation} with the resolvent convention of
  the book.  Re-derive it using $(H_0-z)^{-1}$ and explain the sign change.

  \item Starting from left and right null vectors of
  $A(z)=\identity-K(z)$, verify the Laurent coefficient in
  Eq.~\eqref{eq:fredholm-inverse-residue}.  Show explicitly that it is
  invariant under reciprocal rescaling of the null vectors.

  \item On $L^2([0,1],\rmd E)$, prove that multiplication by a nonzero
  essentially bounded function is not compact.  Reconcile this result with
  compactness of a localized Birman--Schwinger kernel.
\end{enumerate}
